\section{In-Context Gibbs}\label{sec:method}
% Source: G4/METHOD_SPEC.md §2, §4, §5; code: G4/code/icgibbs/sampler.py, g7/methods.py.

\paragraph{The full conditional is one forward pass.} Exchangeability turns a TFM into a Gibbs kernel.
For every $i$ and every assignment $y$ with $p(y_{-i}\mid X)>0$,
\begin{equation}\label{eq:lemma1}
  p(y_i\mid y_{-i},X)=p(y_i\mid x_i,D_{-i}),\qquad D_{-i}=\{(x_j,y_j)\}_{j\neq i}
\end{equation}
(\cref{lem:fullcond}): the conditional of one latent label given all the others is the PPD of a query point $x_i$ with
the rest of the current assignment as its context. Multiplying by the weak-label factor gives the full conditional of
$\pi$ in \eqref{eq:target}, $\pi(y_i=k\mid y_{-i})\propto\tfm(k\mid x_i,D_{-i})\,\lik_i(k)$. A Gibbs sweep over latent
labels therefore needs only forward passes of a frozen TFM, and the weak labels enter only through $\lik$.

\paragraph{Blocks.} One forward pass per label would cost $n$ passes per sweep. We instead partition all $n$ points
at random into $K$ blocks in every sweep and, for each block $B$, query the TFM once with the context $D_{-B}$ (the other
blocks, with their current labels). For PLL, each free $y_i$, $i\in B$, is drawn from
$\propto\tfm(\cdot\mid x_i,D_{-B})\,\lik_i(\cdot)$; blocks include the points whose label is fixed, since blocking only
free points would bias the composition of the context. For LLP, blocks are groups of whole bags (\emph{bag-blocks}), and
each bag's labels are drawn from $\prod_{i\in B_g}\tfm(y_i\mid x_i,D_{-B})$ conditioned on its counts, exactly, by dynamic
programming. An optional Metropolis--Hastings (MH) correction accepts or rejects the whole block, which makes the block update exact
(\cref{prop:blocks}); it costs one forward pass per changed label. Contexts larger than 2{,}000 points are subsampled
uniformly at random.

\paragraph{Warm start and prediction.} The chain starts from a cheap estimate $P_0$ of the labels: for PLL, the TFM fitted
on one candidate per point chosen uniformly at random, cross-fitted over 5 folds; for LLP, linear DLLP
\citep{dllp}. Initial labels are drawn from $P_0$ restricted by $\lik$ (for LLP, conditioned on the bag counts).
After $T$ sweeps, of which the first $T_0$ are discarded, the prediction for a test point is the Rao--Blackwellised
average over $S$ evenly spaced samples,
\begin{equation}\label{eq:predict}
  \hat p(y_*\mid x_*)=\frac1S\sum_{s=1}^S\tfm(y_*\mid x_*,(X,y^{(s)})),
\end{equation}
computed with the TFM's prediction ensemble. This \emph{inductive} predictor is our default; a
\emph{transductive} variant, in which test points join the chain, is evaluated as an ablation.

\begin{algorithm}[t]
\caption{\ours{} (block version; LLP uses bag-blocks). No TFM weight is updated.}
\label{alg:icg}
\begin{algorithmic}[1]
\Require features $X$, weak labels $\W$ with model $\lik$, frozen TFM $\tfm$, test points $X_*$;
         $K$ blocks, $T$ sweeps, burn-in $T_0$, $S$ samples
\State $P_0\gets$ warm start; draw $y\sim P_0\cdot\lik$ \Comment{LLP: conditioned on bag counts}
\For{$t=1,\dots,T$}
  \State partition all points (LLP: all bags) at random into blocks $B_1,\dots,B_K$
  \For{$k=1,\dots,K$}
    \State $Q\gets\tfm(\cdot\mid x_i,D_{-B_k})$ for all $i\in B_k$ \Comment{one forward pass}
    \State PLL: draw $y_i\propto Q_i\,\lik_i$ for free $i\in B_k$; \quad LLP: draw $y_{B_g}\sim\prod_{i\in B_g}Q_i$ given counts $w_g$, for $B_g\subseteq B_k$
    \State \textit{(optional)} accept the new $y_{B_k}$ with the MH probability of \cref{prop:blocks}(a)
  \EndFor
  \State \textbf{if} $t>T_0$ \textbf{then} store $y^{(t)}$
\EndFor
\State \Return $\hat p(y_*\mid x_*)=\frac1S\sum_s\tfm(y_*\mid x_*,(X,y^{(s)}))$ over $S$ evenly spaced stored samples
\end{algorithmic}
\end{algorithm}

\paragraph{Cost and settings.} A run costs $TK+S$ forward passes and no gradient step. All settings were fixed on
development datasets before any evaluation and are shared by every dataset and both studies: $K=10$, $T=30$, $T_0=10$,
$S=10$, i.e.\ 310 forward passes.

\paragraph{Comparators on the same TFM.} To isolate what the sampler adds, we also run with the same TFM: the
\emph{naive TFM}, conditioned on one label per point consistent with the weak labels (a uniformly chosen candidate; a
random count-consistent assignment per bag); \emph{hard-EM}, with the same warm start and blocks but the constrained
arg max instead of sampling (10 iterations); \emph{stochastic EM}, which predicts from the last sample of the chain; and an
\emph{oracle} given the true labels.
